\documentclass[12pt]{article}

\usepackage{amsmath}
\usepackage{amssymb}
\usepackage[a4paper, margin=2.5cm]{geometry}

\title{Equations}
\author{}
\date{}

\begin{document}

\maketitle

\section{Households}

The representative household solves:

\begin{align}
\max_{\{C_t,I_t,K_{t+1}\}_{t=0}^{\infty}}
\sum_{t=0}^{\infty} \beta^t \ln C_t
\end{align}

subject to the budget constraint:

\begin{align}
C_t + I_t + T_t
=
w_t L_t + r_t K_t + \Pi_t,
\end{align}

and the law of motion for capital:

\begin{align}
K_{t+1}
=
(1-\delta)K_t + I_t.
\end{align}

Labor supply is inelastic and fixed:

\begin{align}
L_t = \bar{L}.
\end{align}

The Euler equation is:

\begin{align}
\frac{1}{C_t}
=
\beta
\frac{1}{C_{t+1}}
(1+r_{t+1}-\delta).
\end{align}

In steady state, the Euler equation implies:

\begin{align}
r
=
\frac{1}{\beta}
-
1
+
\delta.
\end{align}

Steady-state investment satisfies:

\begin{align}
I = \delta K.
\end{align}

\section{Final-Good Producer}

The competitive final-good producer combines the outputs of the two intermediate goods sectors into aggregate final output:

\begin{align}
Y_t
=
Y_{1,t}^{\theta}
Y_{2,t}^{1-\theta},
\qquad
0 < \theta < 1.
\end{align}

The final good is the numeraire. The final-good producer solves:

\begin{align}
\max_{Y_{1,t},Y_{2,t}}
\left(
Y_{1,t}^{\theta}Y_{2,t}^{1-\theta}
-
P_{1,t}Y_{1,t}
-
P_{2,t}Y_{2,t}
\right).
\end{align}

The first-order conditions imply the sectoral prices:

\begin{align}
P_{1,t}
&=
\theta
\frac{Y_t}{Y_{1,t}},
\\
P_{2,t}
&=
(1-\theta)
\frac{Y_t}{Y_{2,t}}.
\end{align}

\section{Intermediate Goods Sectors}

There are two intermediate goods sectors, indexed by \(j \in \{1,2\}\). Sector \(j\) produces:

\begin{align}
Y_{j,t}
=
A_j
K_{j,t}^{\alpha_j}
L_{j,t}^{\beta_j}
(\chi_t E_{j,t})^{1-\alpha_j-\beta_j}.
\end{align}

Define the energy share parameter:

\begin{align}
\gamma_j
=
1-\alpha_j-\beta_j.
\end{align}

Equivalently:

\begin{align}
Y_{j,t}
=
A_j
K_{j,t}^{\alpha_j}
L_{j,t}^{\beta_j}
(\chi_t E_{j,t})^{\gamma_j}.
\end{align}

Each intermediate goods firm takes \(P_{j,t}\), \(r_t\), \(w_t\), and \(P_{E,j,t}\) as given and solves:

\begin{align}
\max_{K_{j,t},L_{j,t},E_{j,t}}
\left(
P_{j,t}Y_{j,t}
-
r_t K_{j,t}
-
w_t L_{j,t}
-
P_{E,j,t}E_{j,t}
\right).
\end{align}

The first-order conditions are:

\begin{align}
r_t
&=
P_{j,t}
\alpha_j
\frac{Y_{j,t}}{K_{j,t}},
\\
w_t
&=
P_{j,t}
\beta_j
\frac{Y_{j,t}}{L_{j,t}},
\\
P_{E,j,t}
&=
P_{j,t}
\gamma_j
\frac{Y_{j,t}}{E_{j,t}}.
\end{align}

Under perfect competition, profits in each intermediate goods sector are zero:

\begin{align}
\Pi_{j,t}=0.
\end{align}

Total intermediate-goods-sector profits are:

\begin{align}
\Pi^Y_t
=
\sum_{j=1}^{2}\Pi_{j,t}.
\end{align}

\section{Energy Services}

Each intermediate goods sector combines fossil and renewable energy into energy services:

\begin{align}
E_{j,t}
=
\left[
\omega_j E_{f,j,t}^{\frac{\sigma-1}{\sigma}}
+
(1-\omega_j)E_{r,j,t}^{\frac{\sigma-1}{\sigma}}
\right]^{\frac{\sigma}{\sigma-1}}.
\end{align}

The sector-specific energy price index is:

\begin{align}
P_{E,j,t}
=
\left[
\omega_j^{\sigma}P_{f,t}^{1-\sigma}
+
(1-\omega_j)^{\sigma}P_{r,t}^{1-\sigma}
\right]^{\frac{1}{1-\sigma}}.
\end{align}

Conditional demand for fossil energy is:

\begin{align}
E_{f,j,t}
=
\omega_j^{\sigma}
\left(
\frac{P_{f,t}}{P_{E,j,t}}
\right)^{-\sigma}
E_{j,t}.
\end{align}

Conditional demand for renewable energy is:

\begin{align}
E_{r,j,t}
=
(1-\omega_j)^{\sigma}
\left(
\frac{P_{r,t}}{P_{E,j,t}}
\right)^{-\sigma}
E_{j,t}.
\end{align}

\section{Energy Sector}

The energy sector produces fossil and renewable energy by transforming the aggregate final good into energy inputs.

Fossil energy production is:

\begin{align}
E_{f,t}
=
\frac{z_f}{c_f}M_{f,t}.
\end{align}

Renewable energy production is:

\begin{align}
E_{r,t}
=
\frac{z_r}{c_r}M_{r,t}.
\end{align}

Equivalently, the final-good inputs required for energy production are:

\begin{align}
M_{f,t}
&=
\frac{c_f}{z_f}E_{f,t},
\\
M_{r,t}
&=
\frac{c_r}{z_r}E_{r,t}.
\end{align}

Under perfect competition, fossil and renewable energy prices are equal to marginal costs net of the renewable energy subsidy:

\begin{align}
P_{f,t}
&=
\frac{c_f}{z_f},
\\
P_{r,t}
&=
\frac{c_r}{z_r}
-
s_t.
\end{align}

The subsidy is assumed to satisfy:

\begin{align}
0
\leq
s_t
<
\frac{c_r}{z_r}.
\end{align}

Energy-sector profits are:

\begin{align}
\Pi^E_t
=
P_{f,t}E_{f,t}
+
P_{r,t}E_{r,t}
-
M_{f,t}
-
M_{r,t}
+
s_tE_{r,t}.
\end{align}

Using the pricing equations and perfect competition:

\begin{align}
\Pi^E_t = 0.
\end{align}

Total profits rebated to households are:

\begin{align}
\Pi_t
=
\Pi^Y_t
+
\Pi^E_t.
\end{align}

\section{Government}

The government subsidizes renewable energy production. The subsidy is financed by a lump-sum tax on households:

\begin{align}
T_t
=
s_t E_{r,t}.
\end{align}

\section{Resource Constraint}

Aggregate final output is used for consumption, investment, and as an input in energy production:

\begin{align}
Y_t
=
C_t
+
I_t
+
M_{f,t}
+
M_{r,t}.
\end{align}

Using the linear energy production technologies, this can be written as:

\begin{align}
Y_t
=
C_t
+
I_t
+
\frac{c_f}{z_f}E_{f,t}
+
\frac{c_r}{z_r}E_{r,t}.
\end{align}

\section{Market Clearing}

Capital market clearing is:

\begin{align}
K_t
=
K_{1,t}
+
K_{2,t}.
\end{align}

Labor market clearing is:

\begin{align}
L_t
=
L_{1,t}
+
L_{2,t}.
\end{align}

Total labor supply is fixed:

\begin{align}
L_t = \bar{L}.
\end{align}

Fossil energy market clearing is:

\begin{align}
E_{f,t}
=
E_{f,1,t}
+
E_{f,2,t}.
\end{align}

Renewable energy market clearing is:

\begin{align}
E_{r,t}
=
E_{r,1,t}
+
E_{r,2,t}.
\end{align}

\section{Competitive Equilibrium}

A competitive equilibrium is a sequence of quantities:

\begin{align}
\{
C_t,I_t,K_t,L_t,K_{1,t},K_{2,t},L_{1,t},L_{2,t},
Y_t,Y_{1,t},Y_{2,t},
E_{j,t},E_{f,j,t},E_{r,j,t},
E_{f,t},E_{r,t},M_{f,t},M_{r,t}
\}_{t=0}^{\infty},
\qquad
j \in \{1,2\},
\end{align}

and prices:

\begin{align}
\{
w_t,r_t,P_{1,t},P_{2,t},P_{f,t},P_{r,t},P_{E,j,t}
\}_{t=0}^{\infty},
\qquad
j \in \{1,2\},
\end{align}

such that:

\begin{enumerate}
    \item the representative household satisfies its budget constraint, capital accumulation equation, and Euler equation;
    \item the final-good producer combines the two sectoral goods into aggregate output and satisfies its first-order conditions;
    \item intermediate goods firms maximize profits and satisfy the first-order conditions for capital, labor, and energy demand;
    \item intermediate goods firms minimize the cost of obtaining energy services;
    \item the energy sector satisfies its production and pricing equations;
    \item the government budget constraint holds;
    \item goods, capital, labor, fossil energy, and renewable energy markets clear.
\end{enumerate}

\section{Steady State}

In steady state, all endogenous variables are constant over time. Thus:

\begin{align}
C_t &= C,
&
K_t &= K,
&
I_t &= I,
&
Y_t &= Y,
\\
Y_{j,t} &= Y_j,
&
K_{j,t} &= K_j,
&
L_{j,t} &= L_j,
&
E_{j,t} &= E_j,
\\
E_{f,j,t} &= E_{f,j},
&
E_{r,j,t} &= E_{r,j},
&
E_{f,t} &= E_f,
&
E_{r,t} &= E_r.
\end{align}

Prices are also constant:

\begin{align}
w_t &= w,
&
r_t &= r,
&
P_{1,t} &= P_1,
&
P_{2,t} &= P_2,
\\
P_{f,t} &= P_f,
&
P_{r,t} &= P_r,
&
P_{E,j,t} &= P_{E,j}.
\end{align}

Policy and technology parameters are constant:

\begin{align}
s_t = s,
\qquad
\chi_t = \chi.
\end{align}

The steady-state Euler equation implies:

\begin{align}
r
=
\frac{1}{\beta}
-
1
+
\delta.
\end{align}

Capital accumulation implies:

\begin{align}
I
=
\delta K.
\end{align}

Aggregate final output is produced according to:

\begin{align}
Y
=
Y_1^{\theta}
Y_2^{1-\theta}.
\end{align}

The steady-state resource constraint is:

\begin{align}
Y
=
C
+
\delta K
+
\frac{c_f}{z_f}E_f
+
\frac{c_r}{z_r}E_r.
\end{align}

\section{Additional Variables}

The renewable energy share is defined as:

\begin{align}
\text{Renewable Share}
=
\frac{E_r}{E_f+E_r}.
\end{align}

If the output gap is measured relative to a baseline level of output \(Y^{base}\), then:

\begin{align}
\text{Output Gap}
=
\frac{Y-Y^{base}}{Y^{base}}.
\end{align}

\end{document}